% TOP_PROBLEM: {"id":30007678,"problem_number":"LOCAL-30007678","title":"Vertical integration trade-offs","category_id":21,"category":{"id":21,"name":"economics","display_name":"Economics","description":"Open economic theory statements, published research agendas, and proposed scoped specifications; question provenance and historical priority are qualified per record.","slug":"economics","order_index":21,"created_at":"2026-10-08T00:00:00Z"},"status":"research_question","external_url":null,"legacy_ids":[],"published":false,"statement":"Characterize vertical-integration welfare comparisons when contracts, partners, investment and entry adjust within a specified bargaining model.","economics":{"original_id":167,"field":"Industrial organization and competition","kind":"theoretical","formalization_basis":"adapted_research_question","source_status":"research_agenda","priority_status":"earliest_found","priority_note":"Earliest relevant explicit question or agenda passage located in this audit. No claim of first-ever formulation; earlier versions and later solutions were not exhaustively traced.","earliest_verified_reference":{"authors":["Robin S. Lee","Michael D. Whinston","Ali Yurukoglu"],"title":"Structural Empirical Analysis of Contracting in Vertical Markets","year":2021,"url":"https://robinlee.sites.fas.harvard.edu/papers/HandbookIO_Vol4_VerticalMarkets.pdf","doi":"10.3386/w29282","locator":"Section 5, printed pp. 58–59, future research directions","evidence_summary":"Calls for richer contracts, endogenous partners, incomplete information and dynamic investment."},"current_reference":{"authors":["Robin S. Lee","Michael D. Whinston","Ali Yurukoglu"],"title":"Structural Empirical Analysis of Contracting in Vertical Markets","year":2021,"url":"https://robinlee.sites.fas.harvard.edu/papers/HandbookIO_Vol4_VerticalMarkets.pdf","doi":"10.3386/w29282","locator":"Section 5, printed pp. 58–59, future research directions","evidence_summary":"Calls for richer contracts, endogenous partners, incomplete information and dynamic investment."},"formalization_note":"The source verifies a broader question or research agenda; the population, estimand, equations and restrictions here are our scoped adaptation. This is not a quoted canonical conjecture, and the audit does not prove contemporary unsolved status for every special case. Independent math review tightened feasibility, existence, or estimand assumptions where required."},"statusQualification":"Published question or research agenda verified at the cited locator. The mathematical specification is adapted for this catalogue and includes additional assumptions; source publication does not establish first-ever priority or certify this exact model remains unsolved. The source verifies a broader question or research agenda; the population, estimand, equations and restrictions here are our scoped adaptation. This is not a quoted canonical conjecture, and the audit does not prove contemporary unsolved status for every special case. Independent math review tightened feasibility, existence, or estimand assumptions where required.","statusReviewedAt":"2026-10-08"}
% FIRST_POSTED: 2026-10-08
% ENTRY_KIND: statement_only
% !TeX program = lualatex
\documentclass[11pt]{article}
\usepackage{fontspec}
\usepackage[a4paper,margin=25mm]{geometry}
\usepackage{amsmath,amssymb,mathtools}
\usepackage{xurl}
\usepackage[hidelinks]{hyperref}
\newcommand{\sref}[2]{\hyperlink{source:#1:#2}{[S#2]}}
\setlength{\emergencystretch}{3em}
\begin{document}
% problemId:30007678
\section*{Vertical integration trade-offs}\label{30007678}
\subsection{Definitions and mathematical statement}
Consider an oligopolistic supply chain with a finite upstream-firm set \(U\), finite downstream-firm set \(D\), and consumers with a specified demand system. Firms have positive differentiable production costs and may bargain over contracts combining per-unit prices and fixed transfers. Let \(\theta\) collect demand, costs, bargaining weights, entry costs and the information available during negotiation. In regime \(a=1\), one specified upstream-downstream pair integrates; in regime \(a=0\), all firms remain independent. In either regime firms may change contracting partners, invest and enter, subject to the same technological constraints. Let \(\mathcal E_a(\theta)\) be the set of equilibria under a stated bargaining protocol. Restrict parameter comparisons to cases where both equilibrium sets are nonempty, and define \(W(e)\) as consumer surplus plus firm profits, with profits measured net of real production, investment and entry costs in equilibrium \(e\). Transfers between firms cancel from this welfare measure. Characterize the parameter regions in which every \(e_1\in\mathcal E_1(\theta)\) and every \(e_0\in\mathcal E_0(\theta)\) satisfy \(W(e_1)\geq W(e_0)\). When this strong comparison fails, provide welfare ranges under an explicit equilibrium-selection rule. The task isolates efficiencies, foreclosure and partner rematching within a declared model class. It makes no universal assertion that integration is beneficial or harmful across all possible contracts and information structures.

\par\textbf{Formulation and scope.} The source verifies a broader question or research agenda; the population, estimand, equations and restrictions here are our scoped adaptation. This is not a quoted canonical conjecture, and the audit does not prove contemporary unsolved status for every special case. Independent math review tightened feasibility, existence, or estimand assumptions where required.

\par\textbf{Open-question provenance.} An explicit question or research agenda was verified at the source locator. This does not by itself certify that every later version remains unresolved \sref{30007678}{1}.

\par\textbf{Historical priority.} Earliest relevant explicit question or agenda passage located in this audit. No claim of first-ever formulation; earlier versions and later solutions were not exhaustively traced.

\subsection{Short English statement}
Characterize vertical-integration welfare comparisons when contracts, partners, investment and entry adjust within a specified bargaining model.

\subsection{Sources}
\begin{itemize}\raggedright
\item[\textnormal{[S1]}]\hypertarget{source:30007678:1}{} Robin S. Lee, Michael D. Whinston, Ali Yurukoglu. 2021. Structural Empirical Analysis of Contracting in Vertical Markets. Section 5, printed pp. 58–59, future research directions.
\url{https://robinlee.sites.fas.harvard.edu/papers/HandbookIO_Vol4_VerticalMarkets.pdf}.
DOI: 10.3386/w29282.
{\small\textit{Earliest verified question/agenda reference; historical priority qualified below. Calls for richer contracts, endogenous partners, incomplete information and dynamic investment.}}
\end{itemize}
\end{document}
